\hypertarget{binders-node-0}{%
\chapter{3. Binder's Node 0}\label{binders-node-0}}

Chapter 1 placed the binder among the interposition facilities, and
chapter 2 placed the processes that build a kennel and hold and drop its
privilege. What joins them is the question chapter 2 left standing: how
a confined workload reaches a service it is allowed to use without being
handed a path to it. The answer is the binder bus, and the short form is
that every kennel gets its own and \texttt{kenneld} sits on node 0 of
each as the broker that decides what crosses.

\hypertarget{what-binder-is}{%
\section{3.0 What binder is}\label{what-binder-is}}

A reader who has spent time in Landlock and cgroup BPF may never have
touched binder, because it comes from a corner of the kernel that serves
Android and little else. Binder is Android's inter-process communication
mechanism, and on an Android device it is not one IPC among several but
very nearly the only one: an app calling a system service, an app
handing work to another app, the framework delivering an event, all ride
binder transactions through a kernel driver, brokered by a userspace
registry called the servicemanager.

It belongs to an inverted version of the Unix model. On an ordinary
system a person owns and runs applications and the uid names that
person; Android turns this around and gives each installed app its own
uid, isolating apps from one another by uid the way a multi-user system
isolates people. The users of an Android phone are the apps, mutually
distrusting, and binder is the IPC built for that situation: a
kernel-mediated channel across uid boundaries, with a context manager
that brokers which named service a caller may find and reach, so no app
holds ambient access to another. That shape is the kennel's shape too,
which is the deeper reason the primitive fits rather than merely being
at hand: a kennel is an app-shaped unit of mutual distrust, and the IPC
it needs is the IPC Android already built.

And it is, by a wide margin, the most exercised IPC mechanism in
existence. Binder carries the cross-process traffic of every Android
device, billions of them, and has done so under continuous attack,
fuzzing, and kernel-security scrutiny for more than a decade. Whatever
its quirks, no userspace bus written from scratch will ever see that
volume of adversarial use; taking the driver and its model inherits that
hardening instead of setting out to re-earn it.

\hypertarget{from-connect-to-call}{%
\section{3.1 From connect to call}\label{from-connect-to-call}}

The move the bus makes is from the connection to the call. The AF\_UNIX
shim grants a socket by binding its path into the constructed view and
auditing the workload's connect; that gates once, at connect time, and
says nothing about the authority that rides the protocol afterward. A
D-Bus call can spawn a session process or read a stored secret; a
Wayland socket can read the clipboard. Binder generalises the move the
SSH bastion made first, putting the decision at the operation rather
than the connection: \texttt{kenneld} is the policy decision point for
every protocol call, and the workload holds only unforgeable binder
references, opaque kernel handles with no path to enumerate and no
abstract name to probe. That is the \texttt{reference-monitor} restated
for local services, with the bus as the mediated transaction every call
passes through (\texttt{interpose-as-transaction}).

\hypertarget{the-bus-each-kennel-gets}{%
\section{3.2 The bus each kennel gets}\label{the-bus-each-kennel-gets}}

Each kennel gets its own binderfs instance, a fully independent mount
that shares no nodes with any other kennel's, the way devpts and tmpfs
are independent per mount namespace. binderfs carries
\texttt{FS\_USERNS\_MOUNT}, so the instance mounts inside the kennel's
child user namespace, in the same full-capability context the spawn
already uses to mount tmpfs, devpts, and proc; there is no host-side
mount and no separate privileged step. The spawn mounts a fresh binderfs
at the view's \texttt{/dev/binderfs}, with \texttt{max=256} capping
binder-device allocation per kennel as a denial-of-service bound rather
than a tuning knob, allocates the standard device named \texttt{binder}
through \texttt{BINDER\_CTL\_ADD}, and creates the conventional
\texttt{/dev/binder} symlink so a stock binder client finds its driver
at the default path with no per-workload configuration, the same
principle the socket shim rests on.

Taking node 0 is where the daemon and the instance meet, and the
mechanism is more particular than it looks. \texttt{kenneld} runs in the
initial user namespace and must hold the context-manager fd for every
instance, since one entity owning every node 0 is the precondition for
routing a call from one kennel to another. It cannot be handed that fd:
a binder fd is bound to the process that opened it, so a descriptor
passed over a socketpair cannot be mapped or made context manager by a
different process. \texttt{kenneld} therefore opens the device itself,
through
\texttt{/proc/\textless{}init\textgreater{}/root/dev/binderfs/binder},
which it is permitted to do because the kennel's user namespace is
operator-owned and the operator \texttt{kenneld} is privileged over the
instance. It then calls \texttt{BINDER\_SET\_CONTEXT\_MGR}, which is
one-per-instance, taking node 0. The workload is later granted Landlock
read and write on the device, reached through the \texttt{/dev/binder}
symlink, but never on \texttt{binder-control}: only the spawn setup
allocates devices, and it does so before the seal. When the kennel
exits, the instance disappears with the child's mount namespace, with no
host-side unmount, and pending transactions receive death notifications
as every node is destroyed. Construction is the privileged step here,
not the bus: the \texttt{0\ 0\ 1} map the instance's uid-0-owned nodes
require is the privhelper's act (T5.4), not a binder-specific privileged
surface.

\hypertarget{kenneld-on-node-0}{%
\section{3.3 kenneld on node 0}\label{kenneld-on-node-0}}

Node 0 is the service registry, the analogue of Android's
servicemanager: a caller reaches it through the well-known handle and
never by name, and its verbs carry the same names as Android's interface
even though the transaction codes are \texttt{kenneld}'s own and nothing
is wire-compatible. A service process registers a name with
\texttt{addService}; a workload resolves one with \texttt{getService}
and receives a node reference or a failed reply; \texttt{listServices},
\texttt{isDeclared}, and \texttt{getDeclaredInstances} are the
introspection a caller may run on its own grants. Every name is bounded
and validated, checked against the kennel's settled policy before
\texttt{kenneld} records or resolves it, and audited on every verb. The
check that a service is declared before it can be registered is the
manifest check Android makes against VINTF, with the settled policy
standing in for the manifest: a service the policy does not declare
cannot register and reports as undeclared. Binder's general
capability-passing, the transfer of arbitrary node references between
processes, is not used; references are issued only by \texttt{kenneld},
and they cross between kennels only as \texttt{kenneld} brokers them on
the mesh bus, the subject of the inter-kennel chapter. That drops the
reference-graph bookkeeping that makes Android's implementation subtle
while keeping the unforgeable references, the synchronous calls, the
death notifications, and the per-instance isolation.

The bus carries the kennel's own control plane as well.
\texttt{kennel-bin-init} is a binder consumer on the same instance,
pulling its Plan from node 0 and sending the lifecycle notifications
that tell \texttt{kenneld} a facade crashed or the workload
exec'd.~Those verbs are gated by the unforgeable binder caller identity,
the init's host pid at euid 0, a value the privhelper hands
\texttt{kenneld} at construction rather than anything the caller puts on
the wire, so a workload can address node 0 but cannot exercise them. The
identity binder stamps on every transaction is the gate, not a token the
payload carries, which is why the gate holds against a workload that
controls its own payload bytes entirely.

A reserved prefix, \texttt{org.projectkennel.*}, is \texttt{kenneld}'s
own. It is both context manager and the provider of a built-in set under
that prefix, and two rules hold before any policy lookup: an
\texttt{addService} for a reserved name from any caller other than
\texttt{kenneld} is refused, with no policy override, and a
\texttt{getService} for a reserved name always resolves to
\texttt{kenneld}'s local node and is never routed to a peer kennel. A
reserved node exists only when its policy section is non-empty, so the
absence of the node is itself the proof that the capability was not
granted (\texttt{construction-by-absence}). The registry verbs are cheap
in-memory operations, but a facade verb performs real host I/O, so node
0 is served by a pool of looper threads per instance rather than one:
each looper receives, handles, and replies to its own transactions, and
the pool is sized so that one looper blocked on a host dial does not
stall the others. A dial carries a deadline, so a wedged or hostile
target degrades to a refusal on that one transaction rather than tying
up the instance.

\hypertarget{a-transaction-across-the-boundary}{%
\section{3.4 A transaction across the
boundary}\label{a-transaction-across-the-boundary}}

The AF\_UNIX facade is the example whose whole shape fits here, and it
shows the bus replacing a path with a brokered connection. The workload
sends a \texttt{CONNECT} transaction to the
\texttt{org.projectkennel.IAfUnix/default} node carrying the requested
socket path as a bounded, validated string. \texttt{kenneld} checks the
path against the kennel's \texttt{{[}{[}unix.allow{]}{]}} list, the same
settled grant the shim consumed. On allow it performs the
\texttt{connect()} itself, host-side, and returns the connected fd in
the reply as a \texttt{BINDER\_TYPE\_FD}; on deny it returns a failed
reply and audits it. The workload receives an already-connected socket
and never holds a path into the host's AF\_UNIX namespace; the path does
not appear in the constructed view at all, which closes the residual
where a workload could enumerate the socket paths it was granted, or
probe for ones it was not. The grant is exercised at the call, audited
at the call, and the workload uses the connection without ever reaching
the name behind it (\texttt{mediate-use-not-reach}). The D-Bus facade
and the network facade are larger, each parsing a foreign wire protocol,
and each takes its own chapter; the AF\_UNIX facade is the one that
shows the move whole.

\hypertarget{control-plane-not-data-plane}{%
\section{3.5 Control plane, not data
plane}\label{control-plane-not-data-plane}}

What the workload gets back is an fd, and that is the point at which
\texttt{kenneld} leaves. The connected socket flows its bytes directly
between the workload and the far end; \texttt{kenneld} brokered the
connection and then stepped out of the path it brokered, parsing
control-plane structures and never the payload stream
(\texttt{control-not-data-plane}). The fd and shared-memory object types
are permitted within a single instance, between two parties already
inside the same trust boundary, so passing the connected socket to the
workload is sound.

Reaching a service in another kennel is a separate matter, gated by
bilateral policy: the edge exists only where both kennels' signed policy
declares it, a \texttt{provide} on one side and a \texttt{consume} on
the other (T5.2). How that reach is mediated across the instance
boundary is the subject of the inter-kennel chapter, not this one.

\hypertarget{the-crate-and-the-unsafe}{%
\section{3.6 The crate and the unsafe}\label{the-crate-and-the-unsafe}}

The binder ioctl ABI is hand-rolled in \texttt{kennel-lib-binder}, built
in the same shape as the BPF loader crate, with \texttt{unsafe} confined
to a single file holding the ioctl FFI and the crate listed among the
few that carry it. There is no \texttt{libbinder}: that library and its
NDK form carry Android-specific dependencies, so the crate binds the
kernel's stable binder UAPI directly, the way the BPF crate compiles
against the kernel headers (\texttt{own-your-supply-chain}, and
\texttt{dont-roll-your-own} for what the kernel already provides). The
split between mechanism and policy is the BPF crate's split again:
\texttt{kennel-lib-binder} owns the command loop, the encode and decode
of the command stream and transaction data, the context-manager looper
primitive, device allocation, the protocol-version check that requires
binder version 8 at open, and death-notification plumbing, and
\texttt{kenneld} owns which registrations and resolutions are permitted,
the registries, the reserved services, and cross-kennel mediation, under
\texttt{\#!{[}forbid(unsafe\_code){]}}. The consequence worth stating is
where the untrusted input lands. The command-stream decoder consumes
bytes the workload controls, so it is an untrusted-input parser,
reachable from inside every kennel because the bus is unavoidable, and
it carries a fuzz target for that reason (T5.1). The unsafe is a small,
separately reviewed, fuzzed leaf; the policy brain that decides every
crossing holds none of it (\texttt{rule-of-1},
\texttt{quarantine-the-unsafe}).
